\section{Low rank is enough for the NTK RKHS (three-layer mean kernel)}\label{sec:no-rkhs-advantage}
\noindent
We show that low rank does not shrink the kernel-regime function class: the mean three-layer RF-LR kernel induces the same RKHS as the shallow ReLU kernel.

\subsection{Microscopic distributions: Fisher--Kibble decoupling}\label{subsec:fisher-kibble-main}
We can now state a compact representation under a homogeneity assumption.

\begin{lemma}[Fisher--Kibble decoupling]\label{lem:fisher-kibble}
Let \(x,y\) be input vectors and let \(x_1,y_1\) denote their rank-\(r\) random projections. Define the empirical correlation \( \rho_1 = \cos\angle(x_1,y_1) \) and squared norms \( u=\|x_1\|^2, v=\|y_1\|^2 \). Then:
\begin{itemize}
    \item \( \rho_1 \) follows Fisher's correlation distribution~\cite{fisher1915probable,hotelling1953new}, centered at the true correlation \( \rho \).
    \item \( (u,v) \) follow Kibble's~\cite{kibble1941two} bivariate Gamma (chi-square) law.
    \item Angular and radial parts are independent: \( p(\rho_1,u,v) = p_{\text{Fisher}}(\rho_1)\,p_{\text{Kibble}}(u,v) \).
\end{itemize}
\end{lemma}

Proof: See Appendix~\ref{appendix:fisher-kibble}, Proof~\ref{proof:fisher_kibble}.

\paragraph{Homogeneity assumption enabling decoupling.}
By positive 1-homogeneity and rotational invariance, the base kernel factorizes and \((\rho_1,\|x_1\|,\|y_1\|)\) decouple into angular (Fisher) and radial (Kibble) components; see Appendix~\ref{appendix:homogeneity-decoupling} and Appendix~\ref{appendix:fisher-kibble} for details, explicit densities, and the induced compact three-layer empirical NTK form.

\subsection{Endpoint expansion via hypergeometric analysis}\label{subsec:rkhs-hypergeom}
The RKHS of zonal kernels on the sphere is controlled by the Puiseux behavior near \(\rho=\pm 1\)~\cite{bietti2021deepequalsshallowrelu}. For the mean three-layer kernel, the key step is to analyze the Fisher expectation \(\mathbb{E}[\arccos(\hat{\rho}_r)]\) as \(\rho\to 1\).

\begin{proposition}[Puiseux expansion of \(\mathbb{E}(\arccos(\hat{\rho}_r))\) near \(\rho=1\)]\label{prop:puiseux-arccos-fisher}
Let \(r>2\) and \(\hat{\rho}_r\sim \mathrm{Fisher}(\rho,r)\). Write \(\rho=1-t\) with \(t\downarrow 0\). Then
\[
\mathbb{E}[\arccos(\hat{\rho}_r)]
\;=\;
\sqrt{2t}\,I(r) \;+\; O(t^{3/2}),
\]
where \(I(r)\in(0,\infty)\) admits the closed form
\begin{equation}\label{eq:I-r-closed-form-main}
I(r) \;=\; \frac{(r-2)\,2^{\,r-\frac{5}{2}}}{\sqrt{2\pi}}\;
\frac{\Gamma\!\left(\frac{r-1}{2}\right)\Gamma\!\left(\frac{r}{2}-1\right)}{\Gamma(r-1)}\;
C_1(r),\qquad
C_1(r) \;=\; \frac{\Gamma\!\left(r-\tfrac{1}{2}\right)\Gamma\!\left(r-\tfrac{3}{2}\right)}{\Gamma(r-1)^2}.
\end{equation}
\end{proposition}
\begin{proof}[Proof sketch]
Insert Fisher's density (Appendix~\ref{appendix:homogeneity-decoupling}) into \(\int_{-1}^1 \arccos(s)\,p_{\mathrm{Fisher}}(s\mid 1-t,r)\,ds\) and change variables \(s=1-v\).
Near \(t,v\to 0\), use \(\arccos(1-v)=\sqrt{2v}+O(v^{3/2})\) and apply the hypergeometric connection formula (DLMF §15.8.2) to replace \({}_2F_1(\tfrac12,\tfrac12;r-\tfrac12;1-\tfrac{t+v}{2})\) by the constant \(C_1(r)\) up to \(O((t+v)^{r-3/2})\).
Then the leading integral reduces, after \(w=v/t\), to a Beta integral and yields \eqref{eq:I-r-closed-form-main}. See Appendix~\ref{appendix:proof-conditional-rkhs-homogeneous} for the full proof.
\end{proof}

\begin{proposition}[\(I(r)\) decays as \(1/\sqrt{r}\)]\label{prop:I-r-scaling-main}
Let \(I(r)\) be defined by \eqref{eq:I-r-closed-form-main}. Then \(I(r)=O(1/\sqrt{r})\) as \(r\to\infty\).
\end{proposition}
\begin{proof}[Proof sketch]
Apply Stirling's formula to the Gamma ratios in \eqref{eq:I-r-closed-form-main}. A detailed proof is given in Proposition~\ref{prop:I-r-scaling-appendix} (Appendix~\ref{appendix:I-r-scaling}).
\end{proof}

\subsection{RKHS equivalence via endpoint behavior}\label{subsec:rkhs-bietti-bach}
We now state the RKHS characterization result used to translate Puiseux endpoint behavior into an RKHS statement.

\begin{theorem}[Bietti--Bach RKHS criterion (informal)]\label{thm:bietti-bach-rkhs-criterion}
Let \(K_1,K_2\) be zonal kernels on \(\mathbb{S}^{d-1}\) that are \(C^\infty\) on \([-1,1)\). Suppose that as \(t\downarrow 0\),
\[
K_i(1-t)=K_i(1)-c_i\,t^{\alpha}+O(t^{\alpha+\eta}),\qquad c_i>0,
\]
for the same exponent \(\alpha\in(0,1)\) and some \(\eta>0\) (and likewise at \(\rho=-1\) if needed). Then \(K_1\) and \(K_2\) induce the same RKHS on \(\mathbb{S}^{d-1}\), with equivalent norms.
\end{theorem}
\begin{proof}
This is a direct consequence of~\cite[Theorem~1]{bietti2021deepequalsshallowrelu}, which relates the decay of spherical harmonic coefficients (hence the RKHS) to the Puiseux exponent at the endpoints.
\end{proof}

\begin{corollary}[Low rank is enough for the NTK RKHS: three-layer mean kernel]\label{thm:no_rkhs_advantage}
Under the RF-LR setting with ReLU nonlinearity and isotropic random features on \(\mathbb{S}^{d-1}\), the three-layer mean NTK \(\tilde{\Kop}^{(2)}\) induces the \emph{same} RKHS as the shallow ReLU kernel. In particular, the corresponding RKHSs coincide as sets with equivalent norms.
\end{corollary}
\begin{proof}[Proof sketch]
The mean three-layer kernel is zonal. Proposition~\ref{prop:puiseux-arccos-fisher} gives a Puiseux expansion near \(\rho=1\) with leading exponent \(t^{1/2}\); Proposition~\ref{prop:I-r-scaling-main} shows the \(r\)-dependent coefficient is an \(O(1)\) perturbation of the shallow ReLU coefficient.
Since the shallow ReLU kernel has the same endpoint exponent \(1/2\), Theorem~\ref{thm:bietti-bach-rkhs-criterion} yields RKHS equivalence. See Appendix~\ref{appendix:proof-conditional-rkhs-homogeneous} for a detailed argument.
\end{proof}

\paragraph{Open problem: depth \(L\ge 3\).}
Extending RKHS equivalence to general depth requires controlling endpoint expansions through the full Fisher chain; we leave this to future work (Appendix~\ref{appendix:proof-conditional-rkhs-homogeneous}).

\paragraph{Concentration (appendix).}
We also prove that the empirical three-layer NTK concentrates around its deterministic limit with sub-Gaussian tails in \(r\); see Corollary~\ref{cor:three_layer_ntk_concentration} and Appendix~\ref{appendix:proof-three-layer-ntk-concentration}. The \(I(r)=O(1/\sqrt{r})\) scaling (Proposition~\ref{prop:I-r-scaling-appendix}) is consistent with the fact that Fisher--Kibble corrections to the mean kernel's leading endpoint coefficient vanish as \(r\) grows.
